\section{Experimental Details : Class-Conditional and Unconditional Generation}
\label{sec:pixel_detail}

In this section, we use our safe denoiser in the DMs without text inputs. 
Specifically, we employ experiments on FFHQ~\cite{karras2019style} and ImageNet~\cite{russakovsky2015imagenet} in the $256\times256$ resolution. 
We utilize the pretrained diffusion models from FFHQ~\cite{chung2022diffusion}\footnote{https://github.com/DPS2022/diffusion-posterior-sampling} and ImageNet~\cite{dhariwal2021diffusion}\footnote{https://github.com/openai/guided-diffusion}. 
For the experiments, we use a DDPM solver~\cite{lu2022dpm} with 100 steps.

\paragraph{Unconditional Generation}
For unconditional generation, we utilize the FFHQ dataset to evaluate whether the proposed method effectively mitigates sexual bias, using our method. 
Although FFHQ datset does not include explicit label information, \tabref{ffhq} illustrates that the generated images exibit a noticiable bias toward female images over male ones. 
In this experiment, we select 1K female images from CelebA-HQ~\citep{liu2015faceattributes}\footnote{https://www.kaggle.com/datasets/badasstechie/celebahq-resized-256x256} validation split to serve as unseen negative data, thereby establishing the negative dataset $\{\mathbf{x}^{(1)},...,\mathbf{x}^{(1000)}\}$. 
Then, we employ our safe denoiser to generate 1K images.
While both FFHQ and CelebA-HQ are designed to capture similar distribution, they are not completely aligned. This distinction provides an advantageous experimental setup, where we assess the controllability of image generation using reference images. 
For performance evaluation, we compute FID~\cite{heusel2017gans} score using 1,000 male images from the CelebA-HQ dataset. 
For classification performance, we train a ResNet18 model, as implemented in the PyTorch framework~\footnote{\url{https://pytorch.org/vision/stable/index.html}} using the training dataset in the CelebA-HQ. 
In this experiment, we chose $\sigma=1.0$ and $\eta=0.05$, and employ \textit{Safe denoiser} across the entire denoising timesteps. 

\paragraph{Conditional Generation}
For conditional ImageNet~\cite{russakovsky2015imagenet} experiments at $256 \times 256$ resolution, we use a diffusion model trained on the full ImageNet-256 dataset guided by a classifier~\cite{dhariwal2021diffusion}. 
The diffusion backbone follows a linear noise schedule and is constructed with 1,000 diffusion time-steps. We condition on class labels by scaling the classifier guidance at 5.0, creating a strong pull towards the desired class during the sampling process. 
Each experiment generates 50 samples per class across all 1000 ImageNet classes, producing 50,000 samples that are then evaluated with a pretrained ImageNet classifier for precision, recall, and classification accuracy measurements~\cite{he2016dppresnet50}. 
Our metrics include \textbf{(i) Precision: } the fraction of generated samples that match the designated ImageNet label when conditioned on the class, \textbf{(ii) Recall: } aims to evaluate the diversity and coverage of the targeted class distribution, and \textbf{(iii) Classification Accuracy: } the rate at which generated images are correctly identified as their conditioned label among the 999 classes (excluding the negated target class, i.e, Chihuahua). 
The classification accuracy on the hold-out negated class is also calculated, to evaluate how well the respective method does not generate the negated target class. 
As illustrated in \tabref{imagenet_side_by_side}, we focus on the Chihuahua class to investigate how effectively our proposed safe denoiser can repel a target class while preserving generative quality for other classes in this experiment. 
To avoid unintended Chihuahua generation, aforementioned metrics aim to make sure that samples do not drift toward distinct Chihuahua-like features.
For instance, when we generate an image based on a reference dataset sampled from a Chihuahua, the resulting sample may resemble a Golden Retriever, but it won’t resemble a Chihuahua.

To compare our approach, we implement three variants of the conditional diffusion process: vanilla classifier-guided diffusion model without repellency mechanisms, the Sparse Repellency (SR)~\cite{kirchhof2024sparse} technique applied to the classifier-guided diffusion model, and our safe denoiser technique applied to the same diffusion process. 
For the reference dataset, we select the validation set of Chihuahua class as the negative images.
In this experiment, the safe denoiser technique is applied on the 200 to 800 timesteps of the diffusion process. $\eta=0.02$ was chosen as to control the strength of the repellency away from the Chihuahua target class. 
In the SR variant of the experiment, a repellency scale of 0.01 is combined with a large radius of 300 to push generated samples out of regions resembling the negated target class.